\documentclass[11pt]{article}
\usepackage{icmlko}
\begin{document}

\title{덮음은 공짜가 아니다 --- 축이 늙은 게 아니라 우리가 다른 행에 물어보기 시작한 것}
\author{월드모델 실험실}
\date{노트 493 · 2026년 8월 3일}
\maketitle

\begin{abstract}
노트 492 가 애니의 예외(시계열을 빼면 $+0.0328$)를 ``축이 늙었다''로
닫았다. \textbf{왜 늙었는지}를 파니 늙은 게 아니었다.

\textbf{애니의 트렌드 덮음이 학습 $56.4\%$ 에서 유보 $96.5\%$ 로 뛴다.}
데이터랩 소급 한계($2016$)는 아니다 --- 애니 학습 행 중 $2016$ 이전은
\textbf{$0$ 행}이다. 시기별로 $45.2\% \to 66.1\% \to 55.6\% \to
\mathbf{96.5\%}$ 로, 하필 유보 경계에서 뛴다.

그래서 \textbf{같은 덮음에서} 견줬다. 노트 491 이 찾은 대로 제목이 짧으면
질의가 늘 붙으므로, 제목 길이로 ``늘 붙던 종류''와 ``덮음이 새로 찬 곳''을
가른다.

\begin{center}
\begin{tabular}{lrrr}
\hline
무리 & 학습 & 유보 & 차 \\
\hline
짧은 제목 (늘 붙던 종류) & $0.1914$ & $0.1773$ & $\mathbf{-0.0141}$ \\
긴 제목 (새로 찬 곳) & $0.1266$ & $0.0632$ & $\mathbf{-0.0634}$ \\
\hline
\end{tabular}
\end{center}

\textbf{늘 붙던 행에서는 축이 안 늙었다.} 애니 전체의 늙음
($-0.0592$)에서 짝 맞춘 늙음($-0.0141$)을 빼면 \textbf{$76\%$ 가 선택
이동}이다.
\end{abstract}

\begin{figure}[t]
\centering
\includegraphics[width=\textwidth]{figs/f1.pdf}
\caption{\textbf{가} --- 애니. \textbf{나} --- 같은 덮음에서 견주기.
\textbf{다} --- 세계애니의 같은 무늬.}
\end{figure}

\section{한 번 더 있다}

세계애니에서 같은 일이 $2016$ 경계에 있다 --- 덮음 $12.2\% \to 94.3\%$
그 칸에서 $|rho|$ 가 $0.1874 \to 0.0284$ 로 무너진다(그 뒤 $0.1440 \to
0.1877$ 로 회복).

\textbf{두 사례가 같은 방향이다} --- 덮음이 뛰면 그 칸에서 축이 약해진다.
높은 덮음 자체는 나쁘지 않다(모바일은 내내 $99\%$ 이상인데 $|rho|$ 가
$0.143 \to 0.214$ 로 오른다). \textbf{나쁜 것은 덮음이 \emph{바뀌는}
것이다.}

\section{기계}

질의가 붙는 행은 무작위 표본이 아니다. 제목이 짧고 널리 찾는 것부터
붙는다. 덮음이 $56\%$ 일 때 모형은 \emph{그 $56\%$ 에서} 축을 배운다.
덮음이 $96\%$ 가 되면 나머지 $40\%$ 가 들어오는데, 그 행들은 \textbf{처음
$56\%$ 에 못 들었기 때문에} 축이 약한 행이다. 축은 그대로인데 \textbf{묻는
대상이 바뀌었다.}

\begin{quote}
\textbf{조항} --- \textbf{덮음이 다르면 축 세기를 견주지 않는다.} 학습과
유보의 덮음이 다르면 같은 덮음 아래 부분 무리로 짝을 맞춰 견준다. 안
맞추면 \emph{선택 이동}을 \emph{축의 늙음}으로 읽는다.
\end{quote}

\section{노트 491 에 붙는 단서}

노트 491 이 절단 버그를 고쳐 $+152$ 행을 되찾자고 했다. 이 노트는 거기에
조건을 붙인다.

\begin{quote}
\textbf{새로 붙는 $152$ 행은 지금 붙는 행과 다른 종류다} --- 정의상 지금
질의가 실패하는, 제목이 길고 덜 찾는 쪽이다. 애니 유보에서 그런 행의
$|rho|$ 는 $0.0632$ 로 짧은 쪽($0.1773$)의 $36\%$ 다.
\end{quote}

덮음을 올리는 것은 여전히 옳다 --- 노트 490 이 이 수집에 $-0.0283$ 을
매겼고, 없는 행은 축이 $0$ 이다. 다만 \textbf{$+152$ 행이 판 $+152 \times
(\text{평균 몫})$ 을 준다고 적으면 안 된다.} 되찾은 행의 값은 지금 행의
값보다 작다.

\section*{남은 것}

\begin{itemize}
\item \textbf{애니 덮음이 왜 $2025$ 에 뛰었나} --- 아직 못 댔다. 수집을
언제 돌렸는지 · 그때 어떤 행이 대상이었는지의 기록이 필요하다.
\item \textbf{덮음이 바뀌는 축을 어떻게 다루나} --- 지금 구조는 덮음
표시자를 그냥 준다. 학습 덮음과 유보 덮음이 다르면 그 표시자 자체가
시기 정보다(노트 446 날짜 구멍의 사촌).
\item \textbf{같은 검사를 다른 수집에도} --- 위키 · 태그 · 표지가 시기에
따라 덮음이 바뀌는지 안 봤다.
\end{itemize}

\end{document}
